\section*{Introduction}

Two of the foundational results that have uncovered the relationship between algebraic K-theory and the chromatic filtration are Mitchell's theorem \cite{mitchell} and the Quillen--Lichtenbaum conjecture proved by Voevodsky--Rost \cite{Voevodsky1,Voevodsky2}.
The first says that the algebraic K-theory of any ordinary ring $R$ vanishes at chromatic heights above $1$, namely
\[
	\LTm \KK(R) = 0
	\quad \text{for all $m \geq 2$}.
\]
The second, as reinterpreted by Waldhausen \cite{Waldhausen1984}, says that for suitable rings $R$, the map
\[
	\KK(R)_{(p)} \too \Lof \KK(R),
\]
which approximates K-theory by chromatic heights $0$ and $1$, has bounded above fiber.
We note that the latter implies the former for these suitable rings, since $\Tm$-localization vanishes on bounded above spectra.

Based on these and computational evidence from height $1$, Ausoni--Rognes have outlined the far-reaching redshift conjectures, generalizing these phenomena to higher chromatic heights \cite{RognesOberwolfach,AR02,AR08}.
These conjectures have seen remarkable progress in recent years.
In particular, the higher height analogue of Mitchell's theorem was proved by Clausen--Mathew--Naumann--Noel \cite{DescVan}, using the purity theorem of Land--Mathew--Meier--Tamme \cite{purity}, giving the following redshift upper bound.

\begin{theorem}[{\cite[Theorem C]{DescVan}, \cref{main-thm}}]\label{main-thm-intro}
	If $\CC \in \Catperf$ is $\Lnf$-local then
	\[
		\LTm \KK(\CC) = 0
		\quad \text{for all $m \geq n{+}2$}.
	\]
\end{theorem}

Almost concurrently, Hahn--Wilson proved a higher height analogue of the Quillen--Lichtenbaum conjecture \cite{HW}, later revisited by Angelini-Knoll \cite{AK-BPn}.
They constructed a certain $\bbE_3$-form of the truncated Brown--Peterson spectrum $\BPn$ for which they proved the following.

\begin{thm}[{\cite[Theorem B]{HW}}]\label{QL}
	The following map has a bounded above fiber
	\[
		\KK(\BPn)_{(p)} \too \Lnpf \KK(\BPn).
	\]
\end{thm}

Yuan proved the redshift lower bound for the Lubin--Tate spectrum $\En$ \cite{yuan}.
Burklund--Schlank--Yuan subsequently extended this to all commutative ring spectra using their chromatic nullstellensatz, which expresses one sense in which Lubin--Tate spectra behave like algebraically closed fields \cite{null}.

In this paper, we give a new proof of \cref{main-thm-intro} by descent from $\En$.
The vanishing for the Lubin--Tate spectrum follows from the Quillen--Lichtenbaum result of \cref{QL}, and the passage to the general case relies on the following more familiar sense in which $\En$ is algebraically closed.
The pioneering work of Devinatz--Hopkins \cite{DH}, together with subsequent work of Rognes, Baker--Richter, and Mathew \cite{RognesGalois,BR-Galois,MatGal}, shows that $\En$ is the maximal Galois extension of the $\Kn$-local sphere $\SKn$.
Recent unpublished work of Burklund--Clausen--Levy shows that $\STn$ has the same Galois theory as $\SKn$, and that its algebraic closure coincides with $\En$ as well, which in particular gives the following.

\begin{thm}[{Burklund--Clausen--Levy}]\label{BCL}
	The $\Tn$-local filtered colimit of the $\Tn$-local finite Galois extensions $R$ of $\STn$ is $\En$, that is, there is an isomorphism
	\[
		\LTn(\colim_{R/\STn} R) \iso \En
		\qin \CAlg(\SpTn).
	\]
\end{thm}

\begin{remark}
	Our proof relies on this unpublished result.
	Using instead the $\Kn$-local algebraic closure statement, the argument applies to $\Ln$-local categories.
\end{remark}

We briefly outline the proof.
We start from the vacuous case $n = -1$, where we take $\Lnf$ to be the zero functor, and proceed by induction on the height $n$ as follows:

\begin{steps}
	\item \cref{QL} implies the vanishing for $\En$ via the map from $\BPn$ (\cref{QL-En}).
	\item Since K-theory preserves filtered colimits, and a ring spectrum vanishes if and only if its unit vanishes, \cref{BCL} then gives the vanishing for some finite Galois extension of $\STn$ (\cref{some-vanishing}).
	\item We deduce the vanishing for $\STn$ using a general vanishing descent principle, which may be of independent interest (\cref{descent}, \cref{SpTn-dbl}).
	\item The vanishing for $\STn$ together with the vanishing for $\Lnmf\SS$ from the inductive hypothesis implies the vanishing for $\Lnf\SS$ using the telescopic fracture square (\cref{fracture}, \cref{LnfS}).
	\item An $\Lnf$-local category is a module over $\Perf(\Lnf\SS)$, hence the same is true for their K-theories, so the vanishing of the latter implies the vanishing of the former (\cref{main-thm}).
\end{steps}

A subtlety suppressed in this outline is the distinction between ordinary and $\Tn$-local rings and their module categories.
The telescopic fracture square, together with the inductive hypothesis, shows that this discrepancy is invisible to K-theory localized at the relevant chromatic heights (\cref{fracture}).

\subsection*{Acknowledgements}

The proof in this paper was inspired by Elmanto--Nardin--Yang's proof of Mitchell's theorem \cite{ENY}.
I thank John Rognes and Tomer Schlank for helpful exchanges.
This work was supported by the Max Planck--Weizmann joint postdoctoral program.

I used ChatGPT 5.6 throughout this project, including for locating references and streamlining the exposition, and, notably, for suggesting the proof of \cref{descent}.
